\chapter{Introduction}

The aim of this work is to present a study concerning the development of an
expert system, built in a logic programming environment, for the diagnosis of
anxiety disorders.

Throughout the report we document the various design choices adopted during the
definition phase, as well as the behaviour that the expert system exhibits
during its execution.

Before discussing the project in detail, we provide a brief introduction to the
features that characterise an expert system, to expert systems applied to the
medical field, and to the handling of uncertain knowledge.


\section{Expert Systems}

One of the most actively developed areas in Artificial Intelligence is the study
and design of expert systems, also referred to as \emph{Expert Systems} or
\emph{Knowledge-Based Systems}. In the following sections we illustrate the
features and the components that characterise an expert system.


\subsection{Definition}

The term \emph{expert system} was introduced in 1977 by Feigenbaum, who also
provided a precise definition:

\begin{quotation}
An expert system is a computer program that uses knowledge and reasoning
techniques to solve problems that would normally require the help of a human
expert. An expert system must be able to justify or explain why a particular
solution was reached for a given problem.
\end{quotation}

An expert system must be capable of emulating the work of a human expert in a
given domain. In particular, an expert system must be able to perform the same
actions, provide the same judgements, and offer the same explanations as a
human. This feature, however, should not obscure the ultimate goal of an expert
system. Indeed, contrary to what one might think, its purpose is not to replace
the expert, but rather to assist and help them improve their performance.

The fundamental features of an expert system are the following:

\begin{itemize}
    \item \emph{Possession of specific knowledge}: it is important that an expert
    system possesses enough domain-specific knowledge. The quality of an expert
    system, in fact, does not depend only on the formalism or the inference
    technique adopted, but also on the knowledge that the program possesses;

    \item \emph{Adherence of the domain to reality}: the domain in which an
    expert system operates must not be a simplified model of the real domain, but
    must represent it as it is, of course within certain limits, yet never
    resorting to excessively abstract models;

    \item \emph{Reasoning capability}: an expert system must draw inferences and
    make decisions by exploiting the available knowledge base;

    \item \emph{Capability of interacting with the external world}: an expert
    system must have a communication protocol with the external world. Such
    communication is necessary in order to retrieve the information required to
    carry on the reasoning and to provide explanations to the user;

    \item \emph{Capability of explaining its own operation}: a good expert system
    must be able to make explicit and explain the decisions taken and the
    reasoning by which it reached a conclusion. This capability strengthens the
    degree of trust that the user places in the system;

    \item \emph{Learning capability}: a feature that, while not common to all
    expert systems, is very interesting to consider. Just as a human expert
    learns new experiences and broadens their knowledge over time, an expert
    system should likewise be able to learn new knowledge.
\end{itemize}


\subsection{Components}

Within an expert system it is possible to identify four fundamental components:
the knowledge base, the inference engine, the user interface, and the working
memory.

\subsubsection*{Knowledge Base}

The knowledge base represents the set of information needed to address and solve
a problem in a specific domain.

Such information is called \emph{knowledge}, and it must be represented by means
of a formalism that the machine can interpret (for example, Horn logic,
description logic, modal logic, and so on).

Typically, the knowledge base is organised into two blocks:

\begin{itemize}
    \item \emph{Block of assertions or facts}: it represents the declarative
    knowledge related to a particular problem to be solved, that is, true facts
    introduced at the beginning of the consultation;

    \item \emph{Block of relations or rules}: it contains the relations among the
    elementary facts that are useful for building the body of knowledge. These
    are usually expressed in the form of production rules, or according to other
    types of representation such as semantic networks, frames, scripts, and so
    on.
\end{itemize}

\subsubsection*{Inference Engine}

The inference engine is a program able to deduce, starting from the knowledge
base, the conclusions that constitute the solution of the problem in the domain
in which the expert system operates.

The inference engine is able to identify the set of rules that are applicable at
a given moment, to select the rule to apply among the candidate ones, and to
execute the action indicated by the selected rule.

\subsubsection*{User Interface}

The user interface is responsible for facilitating the interaction between the
user and the expert system. It can be built using user-friendly graphical
interfaces or command-line interfaces. It is important that the interface makes
the dialogue with the user as close as possible to a natural one.

\subsubsection*{Working Memory}

The working memory is the component that allows the storage of newly observed
evidence or of data useful for solving the problem during a specific working
session. The content of this memory will therefore change throughout the entire
execution session of the expert system.


\subsection{Roles}

In general, the roles that interact with an expert system are three:

\begin{itemize}
    \item \emph{Domain expert}: the actor who holds the knowledge and experience
    about a specific domain and is able to solve the problem at hand;

    \item \emph{Knowledge engineer}: the actor who encodes the knowledge,
    provided by the domain expert, into a formalism that the expert system can
    understand;

    \item \emph{User}: the individual who consults the expert system in order to
    benefit from it.
\end{itemize}


\subsection{Building Environments}

The construction of an expert system can be carried out using two different types
of environment:

\begin{itemize}
    \item \emph{Skeleton environments}: these are environments derived from
    existing expert systems, from which the knowledge related to the application
    domain has been removed. Examples include Emycin, Kas, Expert, and so on;

    \item \emph{Environment-type environments}: these are environments that
    provide a language for representing knowledge. Examples include Clips,
    Prolog, and so on.
\end{itemize}

Comparing the two types of environment, one can state that ``Skeleton'' systems
are more immediate, but less flexible, since they are affected by the
specialised architecture of the original product; ``Environment''-type systems,
on the other hand, have flexibility as their strength, while their weakness is
the greater complexity.


\section{Expert Systems in Medicine}

One of the most successful fields in which expert systems have been applied is
the medical one.

Typically, a medical expert system can be used in two different ways: as
decision support or as decision acquisition.

\begin{itemize}
    \item \emph{Decision Support}: expert systems of this type aim to suggest
    options or diagnoses to an expert who bears the responsibility of making
    decisions;

    \item \emph{Decision Acquisition}: expert systems for decision acquisition
    aim to allow a non-expert person to make a decision that goes beyond their
    level of knowledge and experience.
\end{itemize}

The reasons that have driven the spread of expert systems in the medical field
are essentially to be found in the possibility of incorporating, within a
knowledge base, the entire domain-specific experience, isolating it from
possible errors or oversights that may undermine a diagnosis carried out by the
human expert alone. In an expert system, indeed, if a line of reasoning is
correct, it will remain so across subsequent consultations as well.


\section{Handling Uncertain Knowledge}

The cognitive mechanism of the human expert often relies on empirical rules
acquired from their own experience and on data that are not completely certain.

An expert is able to reason even when the available information is uncertain or
incomplete. This implies that conclusions obtained from knowledge held to be
generally true, and by means of approximate reasoning, are not always absolutely
correct.

Starting from these premises, the need emerges for approaches that are able to
replicate this approximate reasoning within expert systems as well.

There are several ways to represent and formalise uncertainty: they range from
quantitative methods (certainty factors, Bayesian networks, belief functions,
the fuzzy approach) to symbolic methods (non-monotonic logics).

All these methods share the common feature of having to deal with the
combination of evidence and the propagation of uncertainty from facts to
conclusions.
